% ============================================================================
%  Why exponential weight averaging (EMA) improves the pruned JORTs subnetwork.
%
%  Standalone-compilable (pdflatex ema_variance_reduction.tex), or strip the
%  preamble and \input the section body into the paper.
%
%  SEAM FOR JORTsE:  the argument below assumes a constant (or slowly varying)
%  effective learning rate eta near the minimizer.  The auto-LR/WD update rule
%  (JORTsE) enters exactly at Eq. (noise-floor): replace the constant eta by the
%  per-epoch effective step eta_t of the automated schedule, and the rewind
%  reset becomes a reset of the accumulators that set eta_t.  Marked <<JORTsE>>.
% ============================================================================
\documentclass[11pt]{article}
\usepackage{amsmath,amssymb,amsthm,bm}
\usepackage[margin=1in]{geometry}

\newtheorem{lemma}{Lemma}
\newtheorem{theorem}{Theorem}
\newtheorem{corollary}{Corollary}
\theoremstyle{remark}
\newtheorem{remark}{Remark}

\newcommand{\R}{\mathbb{R}}
\newcommand{\E}{\mathbb{E}}
\newcommand{\Cov}{\mathrm{Cov}}
\newcommand{\tr}{\mathrm{tr}}
\newcommand{\Hess}{\bm{H}}
\newcommand{\Sig}{\bm{\Sigma}}
\newcommand{\thetabar}{\bar{\bm\theta}}

\begin{document}

\section{Why exponential weight averaging helps}
\label{sec:ema}

We give a mechanism-level argument that replacing the final iterate by an
exponential moving average (EMA) of the weights strictly lowers the expected
loss of the pruned subnetwork \emph{in expectation}, and that the gain is
largest precisely in the regime our recipe creates: a high learning rate
re-injected at the rewind epoch, over a highly sparse (hence ill-conditioned)
model. Throughout, the pruning mask is fixed after the one-shot prune, so
$\bm\theta$ denotes the surviving (unmasked) parameters and $L$ the masked
training loss.

\subsection{Setup: constant-step SGD has an $O(\eta)$ noise floor}

Near a local minimizer $\bm\theta^\star$ of the masked loss, take the quadratic
model
\begin{equation}
L(\bm\theta)\;\approx\;L^\star+\tfrac12(\bm\theta-\bm\theta^\star)^\top
\Hess\,(\bm\theta-\bm\theta^\star),\qquad \Hess\succeq 0 .
\end{equation}
Stochastic gradient descent with step $\eta$ and mini-batch gradient noise of
covariance $\Sig$ (i.e.\ $\nabla L_{\mathrm{batch}}=\Hess(\bm\theta-\bm\theta^\star)+\bm n_t$,
$\E[\bm n_t]=\bm 0$, $\Cov[\bm n_t]=\Sig$) yields, for the centered iterate
$\bm e_t:=\bm\theta_t-\bm\theta^\star$, the first-order vector autoregression
\begin{equation}
\bm e_{t+1}=(\bm I-\eta\Hess)\,\bm e_t-\eta\,\bm n_t .
\label{eq:var1}
\end{equation}
Its stationary covariance $\bm C:=\Cov[\bm e_\infty]$ solves the discrete
Lyapunov equation $\bm C=(\bm I-\eta\Hess)\bm C(\bm I-\eta\Hess)^\top+\eta^2\Sig$,
whose first-order-in-$\eta$ solution is the continuous (Ornstein--Uhlenbeck)
balance
\begin{equation}
\Hess\bm C+\bm C\Hess=\eta\,\Sig
\quad\Longrightarrow\quad
\boxed{\;\bm C=O(\eta)\;}.
\label{eq:noise-floor}
\end{equation}
Thus constant-step SGD does \emph{not} rest at $\bm\theta^\star$; it orbits in a
noise ball whose expected excess loss is the ``noise floor''
\begin{equation}
\E[L(\bm\theta_\infty)]-L^\star=\tfrac12\tr(\Hess\bm C)=O(\eta).
\label{eq:gap-raw}
\end{equation}
% <<JORTsE>>  Replace the constant eta in (noise-floor) by the effective per-epoch
% step eta_t produced by the automated LR/WD rule; C becomes O(eta_t) with eta_t
% set by the JORTsE accumulators, and the rewind resets those accumulators.

\subsection{EMA contracts the noise floor, curvature-weighted}

Let the EMA of the weights with decay $\beta\in(0,1)$ be
$\thetabar_t=\beta\,\thetabar_{t-1}+(1-\beta)\bm\theta_t$
(our runs use $\beta=0.95$), with centered form
$\bar{\bm e}_t=\beta\,\bar{\bm e}_{t-1}+(1-\beta)\bm e_t$.
Work in the eigenbasis of $\Hess$; when $\Hess$ and $\Sig$ are simultaneously
diagonalizable, direction $i$ has curvature $h_i$, stationary iterate variance
$c_i=[\bm C]_{ii}$, and lag-one autocorrelation $\varphi_i=1-\eta h_i\in(0,1)$
from \eqref{eq:var1}.

\begin{lemma}[EMA of an AR(1) sequence]
\label{lem:ar1}
If $x_t$ is stationary with variance $c$ and autocorrelation $\varphi^{|\tau|}$,
and $y_t=\beta y_{t-1}+(1-\beta)x_t$, then
\begin{equation}
\Cov[y]=c\;\underbrace{\frac{1-\beta}{1+\beta}\cdot
\frac{1+\beta\varphi}{1-\beta\varphi}}_{=:~r(\beta,\varphi)} .
\end{equation}
\end{lemma}

\begin{proof}
$y_t=(1-\beta)\sum_{j\ge0}\beta^j x_{t-j}$, so
$\Cov[y]=(1-\beta)^2 c\sum_{j,k\ge0}\beta^{j+k}\varphi^{|j-k|}$. Splitting the
double sum into its diagonal and the two symmetric off-diagonals and summing the
geometric series,
$\sum_{j,k\ge0}\beta^{j+k}\varphi^{|j-k|}
=\tfrac{1}{1-\beta^2}\bigl(1+\tfrac{2\beta\varphi}{1-\beta\varphi}\bigr)
=\tfrac{1}{1-\beta^2}\cdot\tfrac{1+\beta\varphi}{1-\beta\varphi}$.
Multiplying by $(1-\beta)^2$ gives $r(\beta,\varphi)$.
\end{proof}

The contraction factor $r(\beta,\varphi)\le 1$ has two revealing limits:
\begin{equation}
\underbrace{\varphi\to 0}_{\eta h_i\to 1~\text{(stiff)}}:~
r\to\frac{1-\beta}{1+\beta}\ll 1,
\qquad
\underbrace{\varphi\to 1}_{\eta h_i\to 0~\text{(flat)}}:~
r\to 1 .
\label{eq:limits}
\end{equation}
EMA averages away almost all of the noise in the \emph{stiff, high-curvature}
directions and leaves the flat directions untouched --- exactly the right
trade, because the stiff directions are the ones that dominate the loss.

\begin{theorem}[EMA strictly lowers the expected loss]
\label{thm:main}
Under the quadratic--stationary model above, for any $\eta>0$ and
$\beta\in(0,1)$,
\begin{equation}
\E[L(\thetabar_\infty)]-L^\star
=\tfrac12\sum_i h_i\,c_i\,r(\beta,\varphi_i)
\;<\;\tfrac12\sum_i h_i\,c_i
=\E[L(\bm\theta_\infty)]-L^\star ,
\end{equation}
with the gap strictly positive whenever at least one curvature $h_i>0$ (so
$\varphi_i<1$ and $r_i<1$).
\end{theorem}

\begin{proof}
$\E[L(\bm\theta)]-L^\star=\tfrac12\tr(\Hess\Cov[\bm\theta])=\tfrac12\sum_i h_i\,\Cov[\bm\theta_i]$
for any zero-mean stationary $\bm\theta-\bm\theta^\star$. Apply this to the raw
iterate ($\Cov=c_i$) and to the EMA ($\Cov=c_i r_i$ by Lemma~\ref{lem:ar1}). Each
summand obeys $h_i c_i r_i\le h_i c_i$ with strict inequality for $h_i>0$.
\end{proof}

\begin{remark}[General (non-commuting) case]
Without simultaneous diagonalizability the closed form
\eqref{eq:limits} is replaced by a Loewner-order statement: the EMA is a
convex, strictly contractive linear filter of a mean-reverting process, so
$\Cov[\thetabar]\preceq\Cov[\bm\theta]$ and therefore
$\tr(\Hess\,\Cov[\thetabar])\le\tr(\Hess\,\Cov[\bm\theta])$. The direction of
Theorem~\ref{thm:main} survives; only the exact factor changes.
\end{remark}

\subsection{Why the gain is large in \emph{our} recipe}

Theorem~\ref{thm:main} explains that EMA helps; the following corollary explains
why it helped \emph{here}, and why it is complementary to --- not redundant with
--- the hyperparameter rewind.

\begin{corollary}[Rewind $\times$ sparsity amplify the EMA gain]
\label{cor:recipe}
The absolute loss removed by EMA is
\begin{equation}
\Delta:=\tfrac12\sum_i h_i\,c_i\,\bigl(1-r(\beta,\varphi_i)\bigr),
\qquad c_i=O(\eta).
\end{equation}
Hence:
\begin{enumerate}
\item[(a)] \textbf{Rewind.} Resetting the learning rate to a high value at the
prune epoch makes $\eta$ --- and thus every $c_i=O(\eta)$ --- large again, so
$\Delta$ scales with the re-injected step. The rewind supplies the exploration
noise that lets the surviving subnetwork leave the pruning-damaged basin; EMA
harvests the low-loss \emph{mean} of that exploration. This is why the rewind
alone gave no lift (its noise is a cost with no collector), while
rewind$+$EMA does.
\item[(b)] \textbf{Sparsity.} Higher sparsity yields a wider, stiffer Hessian
spectrum (larger $h_{\max}$, more directions with $\eta h_i\!\sim\!1$, i.e.\
$r_i\!\ll\!1$). Both the weights $h_i$ and the contraction $1-r_i$ grow, so
$\Delta$ grows with the target sparsity.
\end{enumerate}
\end{corollary}

\subsection{Instantiation under the automated (JORTsE) schedule}
\label{sec:jortse}

The constant-$\eta$ assumption in \eqref{eq:noise-floor} is the single place where
the automated learning-rate, weight-decay, and momentum rules enter, and making
them explicit turns Corollary~\ref{cor:recipe}(a) from a heuristic into an
identity.

\paragraph{Automated learning rate $=$ self-annealing noise floor.}
The auto-LR rule sets
$\eta_k=\sqrt{b_{k-1}^2+a_{k-1}}-b_{k-1}$ with
$a_{k-1}=\sum_{i<k}\eta_i$ and $b_{k-1}=\sum_{i<k}\eta_i^2$ (Eq.~\eqref{eq:lr} of
the hyperparameter derivation). For $a\ll b^2$ this is
$\eta_k\approx a_{k-1}/(2b_{k-1})$, a monotonically \emph{decaying} step.
Substituting $\eta\to\eta_k$ in \eqref{eq:noise-floor} yields a shrinking noise
floor $c_i=O(\eta_k)$: late in a run the schedule self-anneals and the raw
iterate already tightens, so \emph{EMA's marginal benefit fades as
$\eta_k\to0$}. This is why EMA on a fully annealed run buys little.

\paragraph{Rewind $=$ accumulator reset $=$ noise-floor spike.}
The rewind is precisely a reset of the auto-LR accumulators
$a_{k-1},b_{k-1}\!\to\!0$ at the prune epoch, so $\eta_k$ jumps back to its
warm-up value and $c_i=O(\eta_k)$ is large again. By Corollary~\ref{cor:recipe}
the EMA gain $\Delta$ scales with this re-injected step: the rewind
\emph{manufactures} exactly the high-$\eta$ regime in which EMA is most valuable,
and EMA is what converts the re-injected exploration into a low-loss reported
point. The two are complementary by construction, which is why the rewind alone
gave no lift.

\paragraph{Critically-damped momentum: an exact AR(2) contraction.}
Critical damping sets the per-mode spectral root to
$r_i=1-\sqrt{\eta\lambda_i}\in(0,1)$ (equivalently momentum $r_i^2$): the recurrence
$v^{(k+1)}-(1-\eta\lambda_i+r_i^2)\,v^{(k)}+r_i^2\,v^{(k-1)}=0$ has the repeated root
$r_i$, with $r_s=1-\sqrt{\eta_k\|g_k\|_2^2}$ for salient weights (Hessian eigenvalue
approximated by $\|g_k\|_2^2$) and $r_u=1-\sqrt{\eta_k\alpha_k}$ for unsalient ones.
With stochastic forcing the centered iterate is the critically-damped AR(2) process
of Lemma~\ref{lem:ar2}; stiffness now enters through $\sqrt{\eta\lambda_i}$ rather
than $\eta\lambda_i$. (Here $\beta$ continues to denote the EMA decay of
Lemma~\ref{lem:ar1}; the momentum appears only through its root $r_i$.)

\begin{lemma}[EMA of a critically-damped AR(2)]
\label{lem:ar2}
Let $x_k=2r\,x_{k-1}-r^2\,x_{k-2}+\varepsilon_k$ be the stationary critically-damped
AR(2) process ($0<r<1$, repeated root $r$, white $\varepsilon$), and
$y_k=\beta y_{k-1}+(1-\beta)x_k$ its EMA. Then
\begin{equation}
\frac{\Cov[y]}{\Cov[x]}
=\frac{1-\beta}{1+\beta}\cdot
\frac{\,1+\dfrac{2r\beta(1-r^2)}{1+r^2}-r^2\beta^2\,}{(1-r\beta)^2}
\;=:\;R(\beta,r)\ \le\ 1,
\end{equation}
with lag-one autocorrelation $\rho_1=2r/(1+r^2)$ (Yule--Walker). Its limits recover
Lemma~\ref{lem:ar1} with the root $r$ playing the role of the AR(1) autocorrelation:
$R\to(1-\beta)/(1+\beta)$ as $r\to0$ (stiff, short memory) and $R\to1$ as $r\to1$
(flat, long memory).
\end{lemma}

\begin{proof}
As in Lemma~\ref{lem:ar1}, $y$ is the causal filter
$y_k=(1-\beta)\sum_{j\ge0}\beta^{j}x_{k-j}$, so
$\Cov[y]/\Cov[x]=\tfrac{1-\beta}{1+\beta}\bigl[1+2\sum_{\tau\ge1}\beta^{\tau}\rho_\tau\bigr]$.
The AR(2) autocorrelations satisfy $\rho_\tau=2r\rho_{\tau-1}-r^2\rho_{\tau-2}$
($\rho_0=1$, $\rho_1=2r/(1+r^2)$); their generating function is
$\sum_{\tau\ge0}\beta^{\tau}\rho_\tau=[\,1-2r^3\beta/(1+r^2)\,]/(1-r\beta)^2$, whence
$2\sum_{\tau\ge1}\beta^{\tau}\rho_\tau=[\,4r\beta/(1+r^2)-2r^2\beta^2\,]/(1-r\beta)^2$.
Combining over the common denominator $(1-r\beta)^2$ gives $R(\beta,r)$.
\end{proof}

Because $r_i=1-\sqrt{\eta\lambda_i}<1-\eta\lambda_i$, the momentum process has a
\emph{smaller} spectral root than the overdamped AR(1) iterate of
Lemma~\ref{lem:ar1} at the same $(\eta,\lambda_i)$ --- shorter memory --- so EMA
contracts the stiff modes at least as hard. Since $R(\beta,r_i)\le1$ mode-wise,
Theorem~\ref{thm:main} holds verbatim with $r(\beta,\varphi_i)\!\to\!R(\beta,r_i)$;
the momentum-free Lemma~\ref{lem:ar1} is the $r_i\!\to\!\varphi_i$ special case.

\paragraph{Saliency split (``Jenks before prune, off after'').}
Before the prune epoch the momentum partition is active: salient weights follow
$\beta_s$ (Hessian-curvature damping toward $\bm\theta^\star$) while the
to-be-pruned unsalient weights follow $\beta_u$ (weight-decay damping toward
$\bm 0$), so saliency shapes the training dynamics. After the one-shot prune the
unsalient coordinates are removed and only $\beta_s$ remains --- the saliency
mechanism switches off, and the analysis above (single mean-reverting process on
the surviving weights) applies verbatim to the post-prune, EMA-averaged phase.

\paragraph{Automated weight decay attacks the complementary factor.}
The excess loss is the product $\tfrac12\tr(\Hess\,\bm C)$. The auto-WD rule
$\alpha_k=-P_k+\sqrt{P_k^2+M_k/\eta_k}$ (Eq.~\eqref{eq:wd}) minimizes the running
weight magnitude, targeting the two Zhang \emph{et al.} pruning metrics (small
$\tr\Hess$, small $\sum|\theta|$) and thereby shrinking the $\Hess$ factor; EMA
shrinks the $\bm C$ factor. The two automations act on orthogonal factors of the
same trace, so their gains multiply rather than overlap.

\begin{remark}[Lag-bias in the post-rewind transient, resolved by critical damping]
Lemma~\ref{lem:ar2} bounds the \emph{stationary} variance; immediately after the
rewind the mean has not yet returned to $\bm\theta^\star$. The deterministic
transient of a critically-damped mode decays as $(1+c\,k)\,r_i^{\,k}$, and an EMA
lags a moving mean by ${\approx}\,\beta/(1-\beta)$ steps, so during this transient
the average is \emph{biased} and can carry more error than the raw iterate --- EMA
is not free early. Two facts contain it: (i) the schedule evaluates and checkpoints
the average only after the prune epoch and keeps the best (\texttt{best\_ema}),
discarding the bias-dominated epochs; and (ii) critical damping is the
\emph{shortest} non-oscillatory transient --- a repeated root minimizes settling
time, less damping oscillates and more damping is slower --- so the bias window is
as short as the dynamics allow. The momentum automation thus compresses exactly the
regime in which EMA is unreliable, which is why harvesting the average \emph{after}
the rewind is safe.
\end{remark}

\begin{remark}[Relation to Polyak--Ruppert averaging]
As $\beta\to1$ the EMA becomes a running mean, the forgetting analogue of
Polyak--Ruppert iterate averaging, whose tail average attains the statistically
optimal covariance $\Hess^{-1}\Sig\,\Hess^{-1}$ at the $O(1/t)$ rate,
independent of the (slowly decaying) step. EMA is the non-stationary-friendly
form suited to training that is still moving, e.g.\ just after a rewind.
\end{remark}

\begin{remark}[Assumptions, stated honestly]
The claim is improvement \emph{in expectation}, not for every seed. It uses
(i) a fixed post-prune mask (true: one-shot), (ii) a local quadratic with
stationary gradient noise (accurate as the LR re-anneals after the rewind
spike, not at the spike itself), and (iii) co-diagonalizability of $\Hess,\Sig$
only for the closed form --- the trace inequality holds in general. The
constant-$\eta$ assumption in \eqref{eq:noise-floor} is discharged in
Section~\ref{sec:jortse}, where the automated schedule supplies the explicit
$\eta_k$ and the rewind is realized as the accumulator reset.
\end{remark}

\end{document}
